\SGAchapter{IX}{若干 K 群的计算}{Quelques calculs de groupes K}
\begin{center}P. Berthelot\end{center}
\SourcePage{508}
本讲不对应任何一次口头报告，研讨班后续内容也不使用本讲。这里计算各种\Term{fibration}{纤维化}的若干 $K$ 群：\Term{vector-bundle}{向量丛}、\Term{projective-bundle}{射影丛}、\Term{linear-group}{线性群}下的\Term{torsor}{挠子}等等。本讲始终假设\Term{base-scheme}{基概形}是\Term{noetherian}{诺特的}，记号 $K_\bullet(X)$ 表示\Term{scheme}{概形} $X$ 上\Term{coherent-sheaf}{凝聚层}的\Term{grothendieck-group}{Grothendieck 群}。回忆：若 $X$ 是\Term{noetherian}{诺特的}，则从\Term{coherent}{凝聚} $\mathcal O_X$-\Term{module}{模}的\Term{category}{范畴}的 $K$ 群到 $X$ 上\Term{pseudo-coherent-complex}{伪凝聚复形}的\Term{category}{范畴}的 $K$ 群的\Term{homomorphism}{同态}是\Term{isomorphism}{同构}（第 IV 讲 2.4）。

这里陈述的结果与 Chow 环的相应结果类似，后者可见 Grothendieck 在 1958 年 Chevalley 研讨班上的报告。还可以问：对于按\Term{support}{支集}的\Term{dimension}{维数}定义的\Term{filtration}{滤过}（参见第 X 讲），能在 $K_\bullet$ 的\Term{associated-graded}{相伴分次对象}上得到多大程度的同类命题？这一问题似乎尚无已知答案，甚至命题 1.1（\Term{homotopy}{同伦}\Term{exact-sequence}{正合列}）的情形也是如此，而该命题是下面计算中用得最多的性质之一。

\SGAsection{1}{向量丛}
\begin{SGAresult}{命题}{1.1（\Term{homotopy-exact-sequence}{同伦正合列}）}
设 $X$ 是\Term{noetherian}{诺特}\Term{scheme}{概形}，$X'$ 是 $X$ 的\Term{closed-subscheme}{闭子概形}，$U$ 是 $X$ 中 $X'$ 的补\Term{open-subscheme}{开子概形}，$i$ 是 $X'$ 到 $X$ 的\Term{immersion}{浸入}，$j$ 是 $U$ 到 $X$ 的\Term{immersion}{浸入}。则下列序列\Term{exact-sequence}{正合}：
\SourcePage{509}
\[
K_\bullet(X')\xrightarrow{i_*}K_\bullet(X)\xrightarrow{j^*}K_\bullet(U)\longrightarrow0.
\]
\end{SGAresult}
\begin{proof}
若 $\mathcal F'$ 是\Term{coherent}{凝聚} $\mathcal O_{X'}$-\Term{module}{模}，则 $i_*(\mathcal F')$ 是\Term{support}{支集}包含于 $X'$ 的 $\mathcal O_X$-\Term{module}{模}；因此它在 $U$ 上的\Term{restriction}{限制}为零，并且 $j^*i_*(\operatorname{cl}(\mathcal F'))=0$。由于形如 $\operatorname{cl}(\mathcal F')$ 的元素生成 $K_\bullet(X')$，便有 $j^*\circ i_*=0$。由此得到\Term{homomorphism}{同态}
\[
\rho:K_\bullet(X)/\operatorname{Im}K_\bullet(X')\longrightarrow K_\bullet(U).
\]
现来定义 $\rho$ 的逆\Term{homomorphism}{同态}。设 $\mathcal G$ 是\Term{coherent}{凝聚} $\mathcal O_U$-\Term{module}{模}。由于 $X$ 是\Term{noetherian}{诺特的}，存在\Term{coherent}{凝聚} $\mathcal O_X$-\Term{module}{模} $\mathcal F$，其在 $U$ 上的\Term{restriction}{限制}\Term{isomorphism}{同构}于 $\mathcal G$。若 $\mathcal F'$ 是另一个在 $U$ 上的\Term{restriction}{限制}\Term{isomorphism}{同构}于 $\mathcal G$ 的\Term{coherent}{凝聚} $\mathcal O_X$-\Term{module}{模}，则在 $K_\bullet(X)$ 中有
\[
\operatorname{cl}(\mathcal F)-\operatorname{cl}(\mathcal F')\in\operatorname{Im}K_\bullet(X').
\]
为证明这一点，借助\Term{diagonal-map}{对角映射}把 $\mathcal G$ 看作 $j^*(\mathcal F\oplus\mathcal F')$ 的\Term{subsheaf}{子层}；$\mathcal G$ 于是可以延拓为 $\mathcal F\oplus\mathcal F'$ 的\Term{subsheaf}{子层} $\mathcal F''$。这样就归结为证明
\[
\operatorname{cl}(\mathcal F)-\operatorname{cl}(\mathcal F'')\in\operatorname{Im}K_\bullet(X')
\]
以及对 $\mathcal F'$ 的类似断言。设 $\mathcal H$ 与 $\mathcal K$ 是复合\Term{homomorphism}{同态}
\[
\mathcal F''\longrightarrow\mathcal F\oplus\mathcal F'\longrightarrow\mathcal F
\]
的\Term{kernel}{核}与\Term{cokernel}{余核}。这个\Term{homomorphism}{同态}在 $U$ 上化为 $\mathcal G$ 的\Term{identity-map}{恒等映射}，故 $\mathcal H$ 和 $\mathcal K$ 的\Term{support}{支集}包含于 $X'$。若 $\mathcal I$ 是 $X'$ 在 $X$ 中的\Term{ideal}{理想}，便存在整数 $m,n$ 使得 $\mathcal I^m\mathcal H=0$、$\mathcal I^n\mathcal K=0$（EGA I 9.3.4）。所以可把 $\mathcal H$ 与 $\mathcal K$ 逐层\Term{resolution}{分解}为被 $\mathcal I$ \Term{annihilate}{零化}的\Term{module}{模}，即 $\mathcal O_{X'}$-\Term{module}{模}的\Term{extension}{扩张}，而且有
\SourcePage{510}
\[
\operatorname{cl}(\mathcal H)\in\operatorname{Im}K_\bullet(X'),\qquad
\operatorname{cl}(\mathcal K)\in\operatorname{Im}K_\bullet(X'),
\]
从而得到所需结论。于是，在\Term{coherent}{凝聚} $\mathcal O_U$-\Term{module}{模}的\Term{category}{范畴}上定义一个取值于 $K_\bullet(X)/\operatorname{Im}K_\bullet(X')$ 的函数：对每个\Term{coherent}{凝聚} $\mathcal O_U$-\Term{module}{模}，赋予它在 $X$ 上任一延拓的类。

这个函数具有\Term{additivity}{可加性}。事实上，设
\[
0\longrightarrow\mathcal G'\longrightarrow\mathcal G\longrightarrow\mathcal G''\longrightarrow0
\]
是\Term{coherent}{凝聚} $\mathcal O_U$-\Term{module}{模}的\Term{exact-sequence}{正合列}。设 $\mathcal F$ 是 $\mathcal G$ 在 $X$ 上的延拓；则存在 $\mathcal F$ 的\Term{submodule}{子模} $\mathcal F'$，它延拓 $\mathcal G'$，而 $\mathcal F/\mathcal F'$ 延拓 $\mathcal G''$。所以这个函数经 $K_\bullet(U)$ \Term{resolution}{分解}，定义了所需的\Term{homomorphism}{同态} $K_\bullet(U)\to K_\bullet(X)/\operatorname{Im}K_\bullet(X')$。立即可知它是 $\rho$ 的逆。
\end{proof}

\begin{SGAresult}{命题}{1.2}
设 $X,S$ 是两个\Term{noetherian}{诺特}\Term{scheme}{概形}，$f:X\to S$ 是\Term{flat-morphism}{平坦态射}。给定 $K'(X)$ 中的一族元素 $(\xi_i)_{i\in I}$，考虑\Term{homomorphism}{同态}
\[
\varphi:(K_\bullet(S))^{(I)}\longrightarrow K_\bullet(X),\qquad
\varphi((x_i)_{i\in I})=\sum_{i\in I}\xi_i\cdot f^*(x_i),
\]
其中 $f^*$ 是按第 IV 讲 2.12 定义的从 $K_\bullet(S)$ 到 $K_\bullet(X)$ 的“\Term{inverse-image}{逆像}”\Term{homomorphism}{同态}，并按第 IV 讲 2.10 把 $K_\bullet(X)$ 看作 $K'(X)$-\Term{module}{模}。假设对每个点 $s\in S$，各 $\xi_i$ 在\Term{homomorphism}{同态}
\[
K'(X)\longrightarrow K'(X_s)\longrightarrow K_\bullet(X_s)
\]
下的像生成\Term{abelian-group}{阿贝尔群} $K_\bullet(X_s)$。则 $\varphi$ 是\Term{surjective}{满射}。
\end{SGAresult}
\begin{proof}
\SourcePage{511}
设 $S'$ 是 $S$ 的\Term{closed-subscheme}{闭子概形}，$U$ 是 $S$ 中 $S'$ 的补\Term{open-subscheme}{开子概形}。记 $X_U$ 为 $X$ 在 $U$ 上的\Term{restriction}{限制}，$X_{S'}$ 为 $X$ 在 $S'$ 上的\Term{restriction}{限制}。仍记各 $\xi_i$ 在 $K'(X_{S'})$ 和 $K'(X_U)$ 中的像为 $\xi_i$，并以
\[
\varphi_{S'}:K_\bullet(S')^{(I)}\longrightarrow K_\bullet(X_{S'}),\qquad
\varphi_U:K_\bullet(U)^{(I)}\longrightarrow K_\bullet(X_U)
\]
表示按与 $\varphi$ 相同方式定义的\Term{homomorphism}{同态}。考虑图表
\begin{equation}
\begin{tikzcd}
K_\bullet(S')^{(I)}\arrow[r,"i_*"]\arrow[d,"\varphi_{S'}"']&K_\bullet(S)^{(I)}\arrow[r,"j^*"]\arrow[d,"\varphi"]&K_\bullet(U)^{(I)}\arrow[r]\arrow[d,"\varphi_U"]&0\\
K_\bullet(X_{S'})\arrow[r,"u_*"]&K_\bullet(X)\arrow[r,"v^*"]&K_\bullet(X_U)\arrow[r]&0.
\end{tikzcd}\tag{1.2.1}
\end{equation}
其中 $i,j$ 分别表示 $S',U$ 到 $S$ 的\Term{immersion}{浸入}，$u,v$ 分别表示 $X_{S'},X_U$ 到 $X$ 的\Term{immersion}{浸入}。我们证明：若 $\varphi_{S'}$ 和 $\varphi_U$ 是\Term{surjective}{满射}，则 $\varphi$ 也是\Term{surjective}{满射}。

由 1.1，图中的两横行均\Term{exact-sequence}{正合}。因而，要证明 $\varphi$ 是\Term{surjective}{满射}，只须证明两个方块\Term{commutative-diagram}{交换}。设 $(x_i')_{i\in I}\in K_\bullet(S')^{(I)}$；有
\[
\begin{aligned}
\varphi\circ i_*((x_i'))&=\varphi((i_*(x_i')))=\sum_i\xi_i\cdot f^*(i_*(x_i')),\\
u_*\circ\varphi_{S'}((x_i'))&=u_*\Bigl(\sum_i\xi_i\cdot f_{S'}^*(x_i')\Bigr)
=\sum_i\xi_i\cdot u_*(f_{S'}^*(x_i')),
\end{aligned}
\]
最后一个等号来自\Term{projection-formula}{投影公式}第 IV 讲 2.11.1.2。因此，等式 $\varphi\circ i_*=u_*\circ\varphi_{S'}$ 来自等式 $f^*\circ i_*=u_*\circ f_{S'}^*$；后者由第 IV 讲 3.1.1 成立，因为 $i$ 是\Term{proper-morphism}{固有态射}，$f$ 是\Term{flat-morphism}{平坦态射}。
\SourcePage{512}
设 $(x_i)_{i\in I}\in K_\bullet(S)^{(I)}$；有
\[
\varphi_U\circ j^*((x_i))=\sum_i\xi_i\cdot f_U^*\circ j^*(x_i),
\]
\[
\begin{aligned}
v^*\circ\varphi((x_i))
&=v^*\Bigl(\sum_i\xi_i\cdot f^*(x_i)\Bigr)
=\sum_i\xi_i\cdot v^*\circ f^*(x_i)\\
&=\sum_i\xi_i\cdot f_U^*\circ j^*(x_i),
\end{aligned}
\]
中间的等号来自第 IV 讲 2.12.2。

借助这一结果，可对 $S$ 的\Term{closed-subscheme}{闭子概形}\Term{set}{集合}作\Term{noetherian-induction}{诺特归纳}来证明 1.2。假设对 $S$ 的每个异于 $S$ 的\Term{closed-subscheme}{闭子概形} $S'$，$\varphi_{S'}$ 都是\Term{surjective}{满射}，则上述论证表明：将 $S$ 换为一个\Term{irreducible}{不可约}\Term{affine}{仿射}\Term{open-subscheme}{开子概形}后，可以归结为在\Term{irreducible}{不可约}且\Term{affine}{仿射}的\Term{base-scheme}{基概形}上证明 1.2。此外，图表
\[
\begin{tikzcd}
K_\bullet(S_{\mathrm{red}})^{(I)}\arrow[r]\arrow[d,"\varphi_{S_{\mathrm{red}}}"']&K_\bullet(S)^{(I)}\arrow[d,"\varphi"]\\
K_\bullet(X_{S_{\mathrm{red}}})\arrow[r]&K_\bullet(X)
\end{tikzcd}
\]
表明，还可以假设 $S$ 是\Term{reduced}{约化的}，因为由 1.1，$K_\bullet(X_{S_{\mathrm{red}}})\to K_\bullet(X)$ 是\Term{surjective}{满射}。

现在设\Term{base-scheme}{基概形}是\Term{integral}{整的}且\Term{affine}{仿射的}。令 $S=\operatorname{Spec}(A)$，以 $K$ 表示 $A$ 的\Term{fraction-field}{分式域}；若 $s$ 是 $S$ 的\Term{generic-point}{泛点}，则 $\mathcal O_{S,s}=K$。

设 $y\in K_\bullet(X)$。由假设，$K_\bullet(X_s)$ 由各 $\xi_i$ 的像生成；若 $y_s$ 是 $y$ 在 $K_\bullet(X_s)$ 中的像，则可找到整数 $n_i$，使得
\[
y_s=\sum_i\xi_i\cdot n_i.
\]
\SourcePage{513}
另一方面，有 $K=\varinjlim_f A_f$，其中 $f$ 遍历 $A$ 的元素。因此，$\operatorname{Spec}(K)=\varprojlim_f\operatorname{Spec}(A_f)$，且由 EGA IV 8.2.3 及 8.2.5，
\[
X_s=X\times_S\operatorname{Spec}(K)=\varprojlim_f X\times_S\operatorname{Spec}(A_f).
\]
所考虑的\Term{scheme}{概形}均为\Term{noetherian}{诺特的}，而过渡\Term{morphism}{态射}均为\Term{affine}{仿射}且\Term{flat-morphism}{平坦的}，故由第 IV 讲 3.2.3，
\[
K_\bullet(X_s)=\varinjlim_f K_\bullet(X\times_S\operatorname{Spec}(A_f)).
\]
因为 $K_\bullet$ 上的“\Term{inverse-image}{逆像}”\Term{homomorphism}{同态}关于 $K'$ 是\Term{linear}{线性}的（第 IV 讲 2.12.2），存在 $S$ 的\Term{nonempty}{非空}\Term{open-subscheme}{开子概形} $U=\operatorname{Spec}(A_f)$，使得在 $K_\bullet(X_U)$ 中有
\[
y_U=\sum_i\xi_i\cdot n_i.
\]
设 $S'$ 是\Term{support}{支集}为 $S-U$ 的\Term{closed-subscheme}{闭子概形}。由归纳假设，$\varphi_{S'}$ 是\Term{surjective}{满射}，故图表 (1.2.1) 给出
\[
y-\sum_i\xi_i\cdot n_i\in\operatorname{Im}(\varphi).
\]
由于 $n_i=f^*(n_i)$，便知 $\varphi$ 是\Term{surjective}{满射}。
\end{proof}

\begin{SGAresult}{推论}{1.3}
设 $f:X\to S$ 是\Term{noetherian}{诺特}\Term{scheme}{概形}之间的\Term{flat-morphism}{平坦态射}。假设对 $S$ 的每个点 $s$，\Term{homomorphism}{同态}
\[
K'(X)\longrightarrow K_\bullet(X_s)
\]
是\Term{surjective}{满射}。则由 $f^*$ \Term{derive}{导出}的\Term{homomorphism}{同态}
\[
K_\bullet(S)\otimes_{K'(S)}K'(X)\longrightarrow K_\bullet(X)
\]
是\Term{surjective}{满射}。
\end{SGAresult}
\begin{proof}
\SourcePage{514}
取 $(\xi_i)_{i\in I}$ 为 $K'(X)$ 的所有元素组成的族。这时 1.2 的假设满足，因而\Term{homomorphism}{同态}
\[
\varphi:(K_\bullet(S))^{(K'(X))}\longrightarrow K_\bullet(X)
\]
是\Term{surjective}{满射}。由 $\varphi$ 的定义，这\Term{equivalence}{等价}于\Term{homomorphism}{同态}
\[
K_\bullet(S)\otimes_{K'(S)}K'(X)\longrightarrow K_\bullet(X)
\]
是\Term{surjective}{满射}。
\end{proof}

\begin{SGAresult}{推论}{1.4}
设 $f:X\to S$ 是\Term{noetherian}{诺特}\Term{scheme}{概形}之间的\Term{flat-morphism}{平坦态射}。假设对 $S$ 的每个点 $s$，\Term{canonical}{典范}\Term{homomorphism}{同态} $\mathbb Z\to K_\bullet(X_s)$ 是\Term{isomorphism}{同构}。则\Term{homomorphism}{同态}
\[
f^*:K_\bullet(S)\longrightarrow K_\bullet(X)
\]
是\Term{surjective}{满射}。若 $f$ 还具有\Term{tor-dimension}{Tor 维数}\Term{finite}{有限}的\Term{section}{截面} $g$，则 $f^*$ 是\Term{isomorphism}{同构}。
\end{SGAresult}
\begin{proof}
取仅含元素 $1\in K'(X)$ 的族 $(\xi_i)_{i\in I}$；这时 $\varphi=f^*$，且 1.2 的条件满足，故 $f^*$ 是\Term{surjective}{满射}。

若 $f$ 有\Term{tor-dimension}{Tor 维数}\Term{finite}{有限}的\Term{section}{截面} $g$，便可定义\Term{homomorphism}{同态} $g^*:K_\bullet(X)\to K_\bullet(S)$（第 IV 讲 2.12）。由第 IV 讲 2.12.1.1，$g^*\circ f^*=(f\circ g)^*=\operatorname{Id}$，所以 $f^*$ 是\Term{injective}{单射}，从而是\Term{isomorphism}{同构}。
\end{proof}

\textbf{注记 1.5。} $g$ 的\Term{tor-dimension}{Tor 维数}\Term{finite}{有限}，当且仅当 $X$ 的\Term{fibre}{纤维}在 $g(S)$ 的各点\Term{regular}{正则}，即 $f$ 在 $g(S)$ 的各点是\Term{regular-morphism}{正则态射}（若假设 $f$ 是\Term{finite-presentation}{有限表现的}，也就是说 $f$ 在这些点\Term{smooth}{光滑}）。
\SourcePage{515}
这些性质在 $X$ 与 $S$ 上都是\Term{local-scheme}{局部}的，所以可假设二者都是\Term{affine}{仿射的}；设 $S=\operatorname{Spec}(A)$、$X=\operatorname{Spec}(B)$。这时由 EGA IV 12.3.2，有
\[
\operatorname{tor.dim}_B(A)=\operatorname{dim.proj}_B(A)
=\sup_{s\in S}\operatorname{dim.proj}_{B\otimes_A k(s)}(k(s)),
\]
其中关于 $B$ 的\Term{finite-presentation}{有限表现}假设可以换成\Term{noetherian}{诺特性}假设。而 $k(s)$ \Term{isomorphism}{同构}于 $\mathcal O_{X_s,g(s)}$ 的\Term{residue-field}{剩余域}，且由 EGA $0_{\mathrm{IV}}$ 17.2.5，有
\[
\operatorname{dim.proj}_{B\otimes_A k(s)}(k(s))
=\operatorname{dim.proj}_{\mathcal O_{X_s,g(s)}}(k(s)).
\]
由 EGA $0_{\mathrm{IV}}$ 17.2.11，最后这个数\Term{finite}{有限}，当且仅当 $\mathcal O_{X_s,g(s)}$ 的\Term{cohomological-dimension}{上同调维数}\Term{finite}{有限}，也就是该环\Term{regular}{正则}。

\begin{SGAresult}{命题}{1.6}
设 $S$ 是\Term{noetherian}{诺特}\Term{scheme}{概形}，$\mathcal E$ 是\Term{finite-type}{有限型}\Term{locally-free-module}{局部自由} $\mathcal O_S$-\Term{module}{模}，$X$ 是与 $\mathcal E$ 关联的\Term{vector-bundle}{向量丛}，$f$ 是 $X$ 的\Term{structural-morphism}{结构态射}。则 $f^*:K_\bullet(S)\to K_\bullet(X)$ 是\Term{isomorphism}{同构}。
\end{SGAresult}
\begin{proof}
由于 $f$ \Term{smooth}{光滑}，且有一个\Term{section}{截面}，即\Term{zero-section}{零截面}，要证明 1.6，只须证明对每个 $s\in S$，有 $K_\bullet(X_s)\simeq\mathbb Z$（1.4）。这样就归结到\Term{base-scheme}{基概形}是域的情形。

此时，$\mathcal E$ 是秩为 $1$ 的\Term{submodule}{子模}的\Term{direct-sum}{直和}，因而可把 $f$ \Term{resolution}{分解}为一连串秩为 $1$ 的\Term{fibration}{纤维化}，从而归结到 $X$ 的秩为 $1$，即 $X=\operatorname{Spec}(k[T])$ 的情形。命题于是来自下面的\Term{lemma}{引理}。
\end{proof}

\SourcePage{516}
\begin{SGAresult}{\Term{lemma}{引理}}{1.7}
设 $A$ 是\Term{principal-ideal-domain}{主理想整环}，$X=\operatorname{Spec}(A)$。则\Term{canonical}{典范}\Term{homomorphism}{同态}
\[
\mathbb Z\longrightarrow K_\bullet(X)
\]
是\Term{isomorphism}{同构}。
\end{SGAresult}
\begin{proof}
由于 $X$ \Term{regular}{正则}，\Term{homomorphism}{同态} $K'(X)\to K_\bullet(X)$ 是\Term{isomorphism}{同构}。因此，$K_\bullet(X)$ 由\Term{finite-type}{有限型}\Term{locally-free-module}{局部自由} $\mathcal O_X$-\Term{module}{模}的类生成，也就是由\Term{finite-type}{有限型}\Term{projective-module}{投射} $A$-\Term{module}{模}的类生成。由于 $A$ 是\Term{principal-ideal-domain}{主理想整环}，每个\Term{finite-type}{有限型}\Term{projective-module}{投射模}都是自由的，这就给出\Term{surjective}{满射性}。\Term{injective}{单射性}则来自 $K'(X)$ 是 $\mathbb Z$-\Term{augmented-ring}{增广的}。
\end{proof}

\begin{SGAresult}{命题}{1.8}
设 $S$ 是\Term{noetherian}{诺特}\Term{scheme}{概形}，$\mathcal E$ 是秩为 $r$ 的\Term{locally-free-module}{局部自由} $\mathcal O_S$-\Term{module}{模}，$X$ 是与 $\mathcal E$ 关联的\Term{vector-bundle}{向量丛}，$f$ 是 $X$ 的\Term{structural-morphism}{结构态射}，$g$ 是 $X$ 的\Term{zero-section}{零截面}。则复合\Term{homomorphism}{同态}
\[
K_\bullet(S)\xrightarrow{g_*}K_\bullet(X)\xrightarrow{g^*}K_\bullet(S)
\]
（其中 $g^*$ 按第 IV 讲 2.12.1 和 1.5 定义）是乘以 $K'(S)$ 的元素 $\lambda_{-1}(E)$ 的映射；这里 $E$ 表示 $\mathcal E$ 在 $K'(S)$ 中的类，并置
\[
\lambda_{-1}(E)=\sum_{i=0}^r(-1)^i\lambda^i(E).
\]
\end{SGAresult}
这一命题是第 VII 讲 2.7 对 $K_\bullet$ 的类似命题。
\begin{proof}
设 $\mathcal F$ 是\Term{coherent}{凝聚} $\mathcal O_S$-\Term{module}{模}，其在 $K_\bullet(S)$ 中的类为 $F$。则 $g_*(F)$ 是\Term{complex}{复形} $\mathrm Rg_*(\mathcal F)=g_*(\mathcal F)$ 在 $K_\bullet(X)$ 中的类，也就是通过 $g$ 把 $\mathcal F$ 看作 $\mathcal O_X$-\Term{module}{模}时的类。要计算 $g^*(g_*(F))$，便须取
\[
\mathrm Lg^*(\mathcal F)=\mathcal F\mathbin{\mathop\otimes^{\mathrm L}_{\mathcal O_X}}\mathcal O_S
\]
的类。
\SourcePage{517}
为此，考虑\Term{complex}{复形}
\[
0\longrightarrow\bigwedge^r\mathcal E\otimes_{\mathcal O_S}\mathrm S(\mathcal E)
\longrightarrow\bigwedge^{r-1}\mathcal E\otimes_{\mathcal O_S}\mathrm S(\mathcal E)
\longrightarrow\cdots\longrightarrow\mathcal E\otimes_{\mathcal O_S}\mathrm S(\mathcal E)
\longrightarrow\mathrm S(\mathcal E)\longrightarrow0,
\]
其中 $\mathrm S(\mathcal E)$ 是 $\mathcal E$ 的\Term{symmetric-algebra}{对称代数}。这是 $\mathcal O_S$ 在 $\mathrm S(\mathcal E)$ 上的一个\Term{resolution}{解消}（第 VI 讲 1.11）。记 $X$ 上的关联\Term{complex}{复形}为 $\mathcal L_\bullet$；其第 $k$ 项是 $\bigwedge^k(f^*(\mathcal E))$，所以它是 $\mathcal O_S$ 的一个由\Term{finite-type}{有限型}\Term{locally-free-module}{局部自由} $\mathcal O_X$-\Term{module}{模}组成的\Term{resolution}{解消}；因此有
\[
\mathcal L_\bullet\otimes_{\mathcal O_X}\mathcal F
=\mathcal O_S\mathbin{\mathop\otimes^{\mathrm L}_{\mathcal O_X}}\mathcal F.
\]
于是
\[
\begin{aligned}
g^*(g_*(F))
&=\operatorname{cl}(\mathcal L_\bullet\otimes_{\mathcal O_X}\mathcal F)
=\sum_{i=0}^r(-1)^i\operatorname{cl}(\mathcal L_i\otimes_{\mathcal O_X}\mathcal F)\\
&=\sum_{i=0}^r(-1)^i\operatorname{cl}\bigl(\bigwedge^i\mathcal E\otimes_{\mathcal O_S}\mathcal F\bigr)\\
&=F\cdot\lambda_{-1}(E).
\end{aligned}
\]
由\Term{additivity}{可加性}即得命题。
\end{proof}

\SGAsection{2}{分裂环面下的主丛}
\textbf{2.1。} 设 $S$ 是\Term{noetherian}{诺特}\Term{scheme}{概形}，$T$ 是 $S$ 上的\Term{torus}{环面}。假设 $T$ 是\Term{split-torus}{分裂的}；于是有\Term{isomorphism}{同构} $T\simeq\mathrm D_S(M)$，其中 $M$ 是秩为 $r$ 的\Term{free-abelian-group}{自由阿贝尔群}，也就是说 $T\simeq\mathbb G_{\mathrm m,S}^r$。设 $P$ 是 $T$ 下的\Term{torsor}{挠子}。则 $P$ 在 $S$ 上是\Term{affine}{仿射的}，由一个 $M$ 型\Term{graded-algebra}{分次} $\mathcal O_S$-\Term{algebra}{代数} $\mathcal A$ 定义；若写成 $\mathcal A=\bigoplus_{m\in M}\mathcal A_m$，则每个 $\mathcal A_m$ 都是\Term{invertible-module}{可逆} $\mathcal O_S$-\Term{module}{模}，且对 $m,m'\in M$，由 $\mathcal A$ 的\Term{multiplication}{乘法}定义的\Term{homomorphism}{同态}
\[
\mathcal A_m\otimes_{\mathcal O_S}\mathcal A_{m'}\longrightarrow\mathcal A_{m+m'}
\]
是\Term{isomorphism}{同构}（SGA 3，第 VIII 讲 4.1）。于是，给每个 $m\in M$ 赋予 $\mathcal A_m$ 在 $K'(S)$ 中的类，便定义一个\Term{multiplication}{乘法}\Term{homomorphism}{同态} $\eta:M\to K'(S)$。

\SourcePage{518}
\begin{SGAresult}{命题}{2.2}
采用 2.1 的记号与假设，设 $f:P\to S$ 是 $P$ 的\Term{structural-morphism}{结构态射}。则\Term{homomorphism}{同态}
\[
f^*:K_\bullet(S)\longrightarrow K_\bullet(P)
\]
是\Term{surjective}{满射}，其\Term{kernel}{核}是 $K_\bullet(S)$ 中由各\Term{subgroup}{子群} $(\eta(m)-1)K_\bullet(S)$（$m\in M$）生成的\Term{subgroup}{子群}；\Term{equivalence}{等价}地，只令 $m$ 遍历 $M$ 的一组\Term{generator}{生成元}即可。
\end{SGAresult}
\begin{proof}
先假设 $r=1$，即 $T=\mathbb G_{\mathrm m,S}$。这时 $M=\mathbb Z$，且对每个 $m\in M$，有
\[
\mathcal A_m\simeq\mathcal A_1^{\otimes m}.
\]
因此，若置 $\mathcal A_1=\mathcal L$，则 $\mathcal A=\bigoplus_{n\in\mathbb Z}\mathcal L^{\otimes n}$，这表明 $P$ 等同于\Term{vector-bundle}{向量丛} $X=\mathbf V(\mathcal L)$ 中\Term{zero-section}{零截面}的\Term{complementary-subspace}{补空间}。通过\Term{zero-section}{零截面}把 $S$ 看作 $X$ 的\Term{closed-subscheme}{闭子概形}，由 1.1 得到\Term{exact-sequence}{正合列}
\[
K_\bullet(S)\xrightarrow{g_*}K_\bullet(X)\longrightarrow K_\bullet(P)\longrightarrow0,
\]
其中 $g:S\to X$ 是 $X$ 的\Term{zero-section}{零截面}。另一方面，$f^*$ \Term{resolution}{分解}为
\[
K_\bullet(S)\xrightarrow{f'^*}K_\bullet(X)\longrightarrow K_\bullet(P),
\]
其中 $f'$ 是 $X$ 的\Term{structural-morphism}{结构态射}。由 1.6，$f'^*$ 是\Term{isomorphism}{同构}，其逆为 $g^*$。于是得到\Term{exact-sequence}{正合列}
\[
K_\bullet(S)\xrightarrow{g^*\circ g_*}K_\bullet(S)\xrightarrow{f^*}K_\bullet(P)\longrightarrow0.
\]
由 1.8 得出
\[
K_\bullet(P)\simeq K_\bullet(S)/\lambda_{-1}(L)K_\bullet(S)
=K_\bullet(S)/(1-L)K_\bullet(S),
\]
其中 $L$ 是 $\mathcal L$ 在 $K'(S)$ 中的类。对每个 $n\in\mathbb Z$，$L^n-1$ 都是 $1-L$ 的倍数，这就证明了 $P$ 的秩为 $1$ 时的 2.2。

\SourcePage{519}
接下来对 $P$ 的秩 $r$ 作归纳。假设结论对秩小于 $r$ 的\Term{split-torus}{分裂环面}下的\Term{torsor}{挠子}成立。写 $M=M_1\oplus M_2$，其中 $M_1\simeq\mathbb Z^{r-1}$、$M_2\simeq\mathbb Z$；置
\[
\mathcal A'=\bigoplus_{m\in M_1}\mathcal A_m,\qquad
\mathcal A''=\bigoplus_{m\in M_2}\mathcal A_m,
\]
从而 $\mathcal A=\mathcal A'\otimes_{\mathcal O_S}\mathcal A''$。设 $P_1=\operatorname{Spec}(\mathcal A')$；则 $P_1$ 是 $\mathrm D_S(M_1)$ 下的\Term{torsor}{挠子}，而 $P$ 是 $P_1$ 上群为 $\mathbb G_{\mathrm m,P_1}$ 的\Term{torsor}{挠子}，其\Term{algebra}{代数}是 $\mathcal A''$ 在 $P_1$ 上的\Term{inverse-image}{逆像}。由前面的结果，$K_\bullet(P)$ 是 $K_\bullet(P_1)$ 对由 $(\eta(m_2)-1)K_\bullet(P_1)$（$m_2\in M_2$）生成的\Term{subgroup}{子群}的商。由归纳假设，$K_\bullet(P_1)$ 是 $K_\bullet(S)$ 对由 $(\eta(m_1)-1)K_\bullet(S)$（$m_1\in M_1$）生成的\Term{subgroup}{子群}的商。利用关系
\[
\eta(m_1+m_2)-1=\eta(m_1)\eta(m_2)-1
=\eta(m_1)(\eta(m_2)-1)+\eta(m_1)-1,
\]
可知 $K_\bullet(P)$ 是 $K_\bullet(S)$ 对由 $(\eta(m)-1)K_\bullet(S)$（$m\in M$）生成的\Term{subgroup}{子群}的商；也可以只令 $m$ 遍历 $M$ 的一组\Term{generator}{生成元}。
\end{proof}

\begin{SGAresult}{推论}{2.3}
采用 2.1 的记号与假设，再假设 $S$ \Term{regular}{正则}；则有\Term{canonical}{典范}\Term{isomorphism}{同构}
\[
K'(P)\simeq K'(S)/((\eta(m)-1)_{m\in M}),
\]
其中 $((\eta(m)-1)_{m\in M})$ 是 $K'(S)$ 中由 $\eta(m)-1$ 生成的\Term{ideal}{理想}，这里 $m$ 遍历 $M$ 的一组\Term{generator}{生成元}。
\end{SGAresult}
\begin{proof}
若 $S$ \Term{regular}{正则}，则\Term{canonical}{典范}\Term{homomorphism}{同态} $K'(S)\to K_\bullet(S)$ 是\Term{isomorphism}{同构}（第 IV 讲 2.5）。同样，这时 $P$ 也\Term{regular}{正则}，所以 $K'(P)\to K_\bullet(P)$ 是\Term{isomorphism}{同构}。于是推论立即由 2.2 得出。
\end{proof}

\SourcePage{520}
\SGAsection{3}{射影丛与旗丛}
\begin{SGAresult}{命题}{3.1}
设 $S$ 是\Term{noetherian}{诺特}\Term{scheme}{概形}，$\mathcal E$ 是秩为 $n$ 的\Term{locally-free-module}{局部自由} $\mathcal O_S$-\Term{module}{模}，$\pi=(p_1,\ldots,p_k)$ 是满足 $\sum_i p_i=n$ 的整数列，$X=\mathrm D_\pi(\mathcal E)$ 是 $\mathcal E$ 的 $\pi$ 型\Term{flag-scheme}{旗概形}。则\Term{canonical}{典范}\Term{homomorphism}{同态}
\[
K_\bullet(S)\otimes_{K'(S)}K'(X)\longrightarrow K_\bullet(X)
\]
是\Term{isomorphism}{同构}。
\end{SGAresult}
\begin{proof}
沿用第 VI 讲 4.5 的记号。因此，以 $\mathcal L_i$ 表示 $\mathcal E$ 在 $X$ 上的\Term{inverse-image}{逆像}的\Term{canonical}{典范}商，以 $\mathcal F_i$ 表示 $\mathcal L_i\to\mathcal L_{i-1}$ 的\Term{kernel}{核}（$i=1,\ldots,k$），以 $\lambda_i^{(j)}$ 表示 $\bigwedge^i\mathcal F_j$ 在 $K'(X)$ 中的类。则各 $\lambda_i^{(j)}$ 生成 $K'(S)$-\Term{algebra}{代数} $K'(X)$（第 VI 讲 4.6）。同样，对 $S$ 的每个点 $s$，各 $\lambda_i^{(j)}$ 在 $K'(X_s)$ 中的像生成 $K'(k(s))=\mathbb Z$ 上的\Term{algebra}{代数} $K'(X_s)$。由于 $X_s$ 是域上的\Term{flag-scheme}{旗概形}，它是\Term{regular}{正则}\Term{scheme}{概形}，所以\Term{homomorphism}{同态} $K'(X_s)\to K_\bullet(X_s)$ 是\Term{isomorphism}{同构}；因而\Term{homomorphism}{同态} $K'(X)\to K_\bullet(X_s)$ 是\Term{surjective}{满射}。由 1.3，这就\Term{pushout}{推出}\Term{homomorphism}{同态}
\[
K_\bullet(S)\otimes_{K'(S)}K'(S)\longrightarrow K_\bullet(X)
\]
是\Term{surjective}{满射}。

为证明\Term{injective}{单射性}，可以先归结到\Term{complete-flag-bundle}{完全旗丛} $\mathrm D(\mathcal E)$ 的情形。事实上，假设
\[
K_\bullet(S)\otimes_{K'(S)}K'(\mathrm D(\mathcal E))\longrightarrow K_\bullet(\mathrm D(\mathcal E))
\]
是\Term{injective}{单射}。由于 $K_\bullet(S)\to K_\bullet(\mathrm D(\mathcal E))$ 经 $K_\bullet(X)$ \Term{resolution}{分解}，这个\Term{homomorphism}{同态}\Term{resolution}{分解}为
\SourcePage{521}
\[
\begin{aligned}
K_\bullet(S)\otimes_{K'(S)}K'(X)\otimes_{K'(X)}K'(\mathrm D)
&\longrightarrow K_\bullet(X)\otimes_{K'(X)}K'(\mathrm D)\\
&\longrightarrow K_\bullet(\mathrm D),
\end{aligned}
\]
且第一个\Term{homomorphism}{同态}是\Term{injective}{单射}。由第 VI 讲 4.3、4.6 和 4.7，$K'(\mathrm D)$ 在 $K'(X)$ 上是自由的，因而这\Term{pushout}{推出}
\[
K_\bullet(S)\otimes_{K'(S)}K'(X)\longrightarrow K_\bullet(X)
\]
的\Term{injective}{单射性}。

现在令 $X=\mathrm D(\mathcal E)$。对 $\mathcal E$ 的秩作归纳。$n=1$ 的情形是平凡的，假设该性质对 $n-1$ 已成立。引入 $P=\mathbb P(\mathcal E)$，并令 $\mathcal F$ 为\Term{canonical}{典范}\Term{homomorphism}{同态} $p^*(\mathcal E)\to\mathcal O_P(1)$ 的\Term{kernel}{核}，其中 $p$ 是 $P$ 的\Term{structural-morphism}{结构态射}。则 $X$ 等同于 $P$ 上 $\mathcal F$ 的\Term{flag-scheme}{旗概形}；由归纳假设可见，只须对 $P$ 证明 3.1。结合第 VI 讲 1.1，所要证明的结果可表述如下。
\end{proof}

\begin{SGAresult}{推论}{3.2}
设 $S$ 是\Term{noetherian}{诺特}\Term{scheme}{概形}，$\mathcal E$ 是秩为 $n$ 的\Term{locally-free-module}{局部自由} $\mathcal O_S$-\Term{module}{模}，$X=\mathbb P(\mathcal E)$ 是关联的\Term{projective-bundle}{射影丛}，$f$ 是 $X$ 的\Term{structural-morphism}{结构态射}，$\xi$ 是\Term{sheaf}{层} $\mathcal O_X(1)$ 在 $K'(X)$ 中的类。则\Term{homomorphism}{同态}
\[
\varphi:(K_\bullet(S))^n\longrightarrow K_\bullet(X),\qquad
\varphi((x_i))=\sum_{i=0}^{n-1}\xi^{-i}f^*(x_i)
\]
是\Term{isomorphism}{同构}。
\end{SGAresult}
\begin{proof}
前面已经证明，这个\Term{homomorphism}{同态}是\Term{surjective}{满射}，因为各 $\xi^i$（$i=0,\ldots,-n+1$）组成 $K'(X)$ 在 $K'(S)$ 上的一组基。为证明它是\Term{injective}{单射}，使用\Term{homomorphism}{同态} $f_*$；由于 $f$ 是\Term{projective-morphism}{射影态射}，$f_*$ 有定义。
\SourcePage{522}
由第 IV 讲 2.12.4，有
\[
f_*\Bigl(\sum_{i=0}^{n-1}\xi^{-i}f^*(x_i)\Bigr)
=\sum_{i=0}^{n-1}x_i f_*(\xi^{-i})=x_0,
\]
最后一个等号来自：对 $-n+1\leq i<0$，有 $f_*(\xi^i)=0$（第 VI 讲 2.8）。假设 $\varphi((x_i))=0$。则上式给出 $x_0=0$；对元素 $\xi^i\varphi((x_i))$（$i=1,\ldots,n-1$）作同样的论证，归纳可知所有 $x_i$ 都为零，从而得到\Term{injective}{单射性}。
\end{proof}

\textbf{3.3。} 仍设 $S$ 是\Term{noetherian}{诺特}\Term{scheme}{概形}，$\mathcal E$ 是秩为 $n$ 的\Term{locally-free-module}{局部自由} $\mathcal O_S$-\Term{module}{模}。记 $V=\mathbf V(\mathcal E)$ 为与 $\mathcal E$ 关联的\Term{vector-bundle}{向量丛}（关于 $\mathcal E$ 是\Term{contravariant}{反变}的！）（SGA 3，第 I 讲 4.6），并把它看作 $S$ 上的\Term{group-scheme}{群概形}。设 $P$ 是 $V$ 下的\Term{torsor}{挠子}；我们要计算 $K_\bullet(P)$。

先回忆如何把 $P$ \Term{embedding}{嵌入} $S$ 上的一个\Term{projective-bundle}{射影丛}。$V$ 下的\Term{torsor}{挠子}对应于 $\mathcal E$ 被 $\mathcal O_S$ 的\Term{extension}{扩张} $\mathcal F$。事实上，若给定\Term{extension}{扩张}
\begin{equation}
0\longrightarrow\mathcal O_S\xrightarrow{s}\mathcal F\xrightarrow{u}\mathcal E\longrightarrow0,
\tag{3.3.1}
\end{equation}
则考虑由 $s$ 定义的\Term{homomorphism}{同态} $\mathbf V(\mathcal F)\to\mathbf V(\mathcal O_S)$，以及 $\mathbf V(\mathcal O_S)=\operatorname{Spec}(\mathcal O_S[T])$ 由 $T=1$ 定义的\Term{section}{截面}，并取
\[
P=\mathbf V(\mathcal F)\times_{\mathbf V(\mathcal O_S)}S,
\]
便得到 $V$ 下的一个\Term{torsor}{挠子}；$V$ 在 $P$ 上的\Term{action}{作用}由它在 $\mathbf V(\mathcal F)$ 上的\Term{action}{作用}诱导。

反过来，设 $P$ 是 $V$ 下的\Term{torsor}{挠子}。以 $\mathcal P$ 表示从 $S$ 到 $P$ 的\Term{homomorphism}{同态}的芽所成的\Term{sheaf}{层}；它是从 $S$ 到 $V$ 的\Term{homomorphism}{同态}的芽所成的\Term{sheaf}{层}下的\Term{principal-homogeneous-sheaf}{主齐性层}，后一个\Term{sheaf}{层}等同于 $\mathcal E$ 的\Term{dual}{对偶} $\mathcal E^\vee$。由此按经典方法得到 $\mathcal O_S$ 被 $\mathcal E^\vee$ 的一个\Term{extension}{扩张}，再取\Term{dual}{对偶}便得到序列 (3.3.1)。

\SourcePage{523}
因此，设 $P$ 是 $V$ 下的\Term{torsor}{挠子}，对应于 $\mathcal E$ 被 $\mathcal O_S$ 的\Term{extension}{扩张} $\mathcal F$。$P$ 的 $\mathcal O_S$-\Term{algebra}{代数} $\mathcal A(P)$ 是 $\mathrm S(\mathcal F)\otimes_{\mathrm S(\mathcal O_S)}\mathcal O_S$，也就是 $\mathrm S(\mathcal F)/(s-1)\mathrm S(\mathcal F)$，它等同于零次\Term{homogeneous-fraction}{齐次分式}所成的环 $\mathrm S(\mathcal F)_{(s)}$。因此，$P$ 等同于 $\mathbb P(\mathcal F)$ 的\Term{open-subscheme}{开子概形} $D_+(s)$。$P$ 在 $\mathbb P(\mathcal F)$ 中的\Term{complementary-subspace}{补空间}于是为 $\operatorname{Proj}(\mathrm S(\mathcal F)/s\cdot\mathrm S(\mathcal F))$，由 (3.3.1)，它等同于 $\mathbb P(\mathcal E)$。所以有\Term{immersion}{浸入}
\[
\begin{tikzcd}
\mathbb P(\mathcal E)\arrow[r,hook]\arrow[dr]&\mathbb P(\mathcal F)\arrow[d]&P\arrow[l,hook']\arrow[dl]\\
&S&
\end{tikzcd}
\]
其中 $\mathbb P(\mathcal E)$ 在 $\mathbb P(\mathcal F)$ 中是闭的，$P$ 是\Term{open-subset}{开集}且为 $\mathbb P(\mathcal E)$ 的补集。

\begin{SGAresult}{命题}{3.4}
设 $S$ 是\Term{noetherian}{诺特}\Term{scheme}{概形}，$\mathcal E$ 是秩为 $n$ 的\Term{locally-free-module}{局部自由} $\mathcal O_S$-\Term{module}{模}，$V=\mathbf V(\mathcal E)$ 是关联的\Term{vector-bundle}{向量丛}，$P$ 是 $V$ 下的\Term{torsor}{挠子}，$f$ 是 $P$ 的\Term{structural-morphism}{结构态射}。则\Term{canonical}{典范}\Term{homomorphism}{同态}
\[
f^*:K_\bullet(S)\longrightarrow K_\bullet(P)
\]
是\Term{isomorphism}{同构}。
\end{SGAresult}
\begin{proof}
采用 3.3 中的等同，有\Term{exact-sequence}{正合列}
\begin{equation}
K_\bullet(\mathbb P(\mathcal E))\xrightarrow{i_*}K_\bullet(\mathbb P(\mathcal F))
\longrightarrow K_\bullet(P)\longrightarrow0,
\tag{3.4.1}
\end{equation}
其中 $i$ 表示由 $u$ 定义的\Term{closed-immersion}{闭浸入} $\mathbb P(\mathcal E)\to\mathbb P(\mathcal F)$。由 3.2，有
\[
\begin{aligned}
K_\bullet(\mathbb P(\mathcal E))&=K_\bullet(S)\otimes_{K'(S)}K'(\mathbb P(\mathcal E)),\\
K_\bullet(\mathbb P(\mathcal F))&=K_\bullet(S)\otimes_{K'(S)}K'(\mathbb P(\mathcal F)).
\end{aligned}
\]
由于 $\mathbb P(\mathcal E)$ 和 $\mathbb P(\mathcal F)$ \Term{smooth}{光滑}，$i$ 是\Term{regular-immersion}{正则浸入}，故有\Term{homomorphism}{同态} $i_*:K'(\mathbb P(\mathcal E))\to K'(\mathbb P(\mathcal F))$；$K_\bullet$ 之间的\Term{homomorphism}{同态}便由 $K'$ 之间的\Term{homomorphism}{同态}与 $K_\bullet(S)$ 作\Term{tensor-product}{张量积}得到。要说明这一点，只须证明：对 $a\in K_\bullet(S)$ 和 $b\in K'(\mathbb P(\mathcal E))$，有
\SourcePage{524}
\[
i_*(b\cdot h^*(a))=i_*(b)\cdot g^*(a),
\]
其中 $g,h$ 是 $\mathbb P(\mathcal F)$ 和 $\mathbb P(\mathcal E)$ 的\Term{structural-morphism}{结构态射}。由于 $h=g\circ i$，由第 IV 讲 2.12.4，有
\[
i_*(b\cdot h^*(a))=i_*(b\cdot i^*g^*(a))=g^*(a)\cdot i_*(b).
\]
因此，由 (3.4.1) 得到
\[
K_\bullet(P)\simeq K_\bullet(S)\otimes_{K'(S)}
\bigl(K'(\mathbb P(\mathcal F))/\operatorname{Im}K'(\mathbb P(\mathcal E))\bigr).
\]
设 $\xi_F,\xi_E$ 是 $\mathbb P(\mathcal F),\mathbb P(\mathcal E)$ 上\Term{canonical}{典范}\Term{invertible-sheaf}{可逆层}的类。由第 VI 讲 1.1，$K'(\mathbb P(\mathcal F))$（相应地，$K'(\mathbb P(\mathcal E))$）在 $K'(S)$ 上有基 $\xi_F^k$，其中 $k=0,\ldots,-n$（相应地，$\xi_E^k$，其中 $k=0,\ldots,-n+1$）。\Term{projection-formula}{投影公式}给出
\[
i_*(\xi_E^k)=i_*(i^*(\xi_F^k))=\xi_F^k i_*(1).
\]
若 $\mathcal I$ 是 $\mathbb P(\mathcal E)$ 在 $\mathbb P(\mathcal F)$ 中的\Term{ideal}{理想}，则 $\mathcal I$ 是\Term{invertible-ideal}{可逆理想}，且 $i_*(1)=1-\mathcal I$。此外，$\mathcal I$ 由 $\mathrm S(\mathcal F)$ 的\Term{ideal}{理想} $J$ 定义，$J$ 是 $\mathrm S(\mathcal F)\to\mathrm S(\mathcal E)$ 的\Term{kernel}{核}，并由 $\mathcal F$ 中\Term{isomorphism}{同构}于 $\mathcal O_S$ 的一个\Term{submodule}{子模}生成；因此，$\mathcal I\simeq\mathcal O_{\mathbb P(\mathcal F)}(-1)$。所以，$K'(\mathbb P(\mathcal E))$ 的像是由 $\xi_F^k-\xi_F^{k-1}$（$k=0,\ldots,-n+1$）生成的 $K'(S)$-\Term{submodule}{子模}。因而\Term{homomorphism}{同态}
\[
K'(S)\longrightarrow K'(\mathbb P(\mathcal F))/\operatorname{Im}(K'(\mathbb P(\mathcal E)))
\]
是\Term{isomorphism}{同构}，\Term{homomorphism}{同态} $K_\bullet(S)\to K_\bullet(P)$ 也就是\Term{isomorphism}{同构}。
\end{proof}

\SGAsection{4}{$\mathrm{GL}(n)_S$ 下的主丛}
\begin{SGAresult}{命题}{4.1}
设 $S$ 是\Term{noetherian}{诺特}\Term{scheme}{概形}，$G$ 是 $S$-群 $\mathrm{GL}(n)_S$，$P$ 是 $G$ 下的\Term{torsor}{挠子}，$\mathcal E$ 是与 $P$ 对应的秩为 $n$ 的\Term{locally-free-module}{局部自由} $\mathcal O_S$-\Term{module}{模}。以 $E$ 表示 $\mathcal E$ 在 $K'(S)$ 中的类，以 $J$ 表示由 $E-n$ 生成的\Term{lambda-ideal}{$\lambda$-理想}。
\SourcePage{525}
若 $f$ 是 $P$ 的\Term{structural-morphism}{结构态射}，则 $f^*$（由于 $f$ \Term{flat-morphism}{平坦}而有定义）是\Term{surjective}{满射}，并定义\Term{isomorphism}{同构}
\[
K_\bullet(S)/J\cdot K_\bullet(S)\xrightarrow{\sim}K_\bullet(P).
\]
\end{SGAresult}
\begin{proof}
设 $\mathcal E_0=(\mathcal O_S)^n$，$e_1,\ldots,e_n$ 为 $\mathcal E_0$ 的显然的一组基，$\mathcal E_0^{(i)}$ 为 $\mathcal E_0$ 中由 $e_1,\ldots,e_i$ 生成的\Term{submodule}{子模}。$G$ 的如下\Term{subfunctor}{子函子}可由 $G$ 的一个\Term{subgroup-scheme}{子群概形} $B$ 表示：它给每个 $S$-\Term{scheme}{概形} $S'$ 赋予 $\mathcal E_{0S'}$ 的所有保持各\Term{submodule}{子模} $\mathcal E_{0S'}^{(i)}$ \Term{invariant}{不变}的\Term{automorphism}{自同构}构成的\Term{subgroup}{子群}。$B$ 是 $G$ 的一个\Term{borel-subgroup}{Borel 子群}。\Term{quotient-sheaf}{商层} $G/B$（关于 fpqc \Term{topology}{拓扑}）可由 $\mathcal E_0$ 的\Term{flag-scheme}{旗概形}表示，记为 $X_0$。设 $\mathcal L_0^{(i)}$ 是 $\mathcal E_0$ 在 $X_0$ 上的\Term{inverse-image}{逆像}的\Term{canonical}{典范}商，$\mathcal G_0^{(i)}$ 是 $\mathcal E_{0X_0}\to\mathcal L_0^{(i)}$ 的\Term{kernel}{核}，并令
\[
\mathcal F_0^{(i)}=\mathcal G_0^{(i-1)}/\mathcal G_0^{(i)}
\simeq\operatorname{Ker}(\mathcal L_0^{(i)}\longrightarrow\mathcal L_0^{(i-1)}).
\]
另一方面，$B$ 是其\Term{unipotent-radical}{幺幂根基} $U$ 与其保持各 $e_i$ \Term{invariant}{不变}的\Term{maximal-torus}{极大子环面} $T$ 的\Term{semidirect-product}{半直积}。设 $Y_0=G/U$。则 $Y_0$ 是以 $X_0$ 为底、群为 $T_{X_0}$ 的\Term{torsor}{挠子}。若 $S'$ 是 $S$-\Term{scheme}{概形}，给定 $Y_0$ 的一个 $S'$-点，\Term{equivalence}{等价}于在 $S'$ 上某个\Term{faithfully-flat}{忠实平坦}且\Term{quasi-compact}{拟紧}的 $S''$ 上给定 $\mathcal E_{0S''}$ 的一个\Term{automorphism}{自同构}，模掉 $U_{S''}$ 中的\Term{unipotent}{幺幂}元素；也就是说，给定 $\mathcal E_{0S''}$ 的一个\Term{flag}{旗}，其中各\Term{submodule}{子模}都是自由的，并且对每个 $i$，给定一类在 $\mathcal O_{S''}$ 上自由的元素 $f_i^\alpha$，满足
\[
f_i^\alpha-f_i^\beta\in\mathcal G_{0S''}^{(i+1)};
\]
也就是给定 $\mathcal E_{0S''}$ 的一个\Term{flag}{旗}，以及对每个 $i$ 给定 $\mathcal F_{0S''}^{(i)}$ 的一组基。这些数据须满足 $S''\times_{S'}S''$ 上的\Term{gluing}{粘合}条件；由 fpqc \Term{descent}{下降}，这\Term{equivalence}{等价}于在 $S'$ 上给定同样的数据。于是立即看出，$Y_0$ 是由如下\Term{algebra}{代数}定义的\Term{affine}{仿射} $X_0$-\Term{scheme}{概形}：
\SourcePage{526}
\[
\mathcal A_0=\coprod_{k_i\in\mathbb Z}
(\mathcal F_0^{(1)})^{\otimes k_1}\otimes
(\mathcal F_0^{(2)})^{\otimes k_2}\otimes\cdots\otimes
(\mathcal F_0^{(n)})^{\otimes k_n}.
\]
群 $G$ \Term{action}{作用}于 $\mathcal E_0$、配备了各\Term{sheaf}{层} $\mathcal L_0^{(i)}$ 的 $X_0$，以及 $Y_0$。把这些\Term{object}{对象}与 $P$ 关于 $G$ 作乘积，便得到与 $P$ 关联的\Term{sheaf}{层} $\mathcal E$、$\mathcal E$ 的\Term{flag-scheme}{旗概形} $X=\mathrm D(\mathcal E)=P/B$，以及\Term{scheme}{概形} $Y=P/U$。$X$ 上的\Term{canonical}{典范}\Term{sheaf}{层}，用与 $X_0$ 上相应\Term{object}{对象}相同的字母表示。于是有图表
\[
\begin{tikzcd}
&Y=P/U\arrow[dr]&\\
P\arrow[ur]\arrow[dr]&&X=P/B=\mathrm D(\mathcal E)\arrow[dl]\\
&S&
\end{tikzcd}
\]
注意，$P$ 是以 $Y$ 为底、\Term{structural-group}{结构群}为 $U_Y$ 的\Term{torsor}{挠子}，$Y$ 是以 $X$ 为底、\Term{structural-group}{结构群}为 $T_X$ 的\Term{torsor}{挠子}，它由如下 $\mathcal O_X$-\Term{algebra}{代数}定义：
\[
\mathcal A=\coprod_{k_i\in\mathbb Z}
(\mathcal F^{(1)})^{\otimes k_1}\otimes\cdots\otimes(\mathcal F^{(n)})^{\otimes k_n}.
\]
因此，\Term{homomorphism}{同态} $K_\bullet(S)\to K_\bullet(P)$ \Term{resolution}{分解}为
\[
K_\bullet(S)\longrightarrow K_\bullet(X)\longrightarrow K_\bullet(Y)\longrightarrow K_\bullet(P).
\]
由 3.1，有 $K_\bullet(X)=K_\bullet(S)\otimes_{K'(S)}K'(X)$。设 $\xi_i$ 是 $\mathcal F^{(i)}$ 在 $K'(X)$ 中的类。由 2.2，$K_\bullet(Y)$ 是 $K_\bullet(X)$ 对由各 $(\xi_i-1)K_\bullet(X)$ 生成的\Term{submodule}{子模}的商。若 $I$ 是 $K'(X)$ 中由各 $\xi_i-1$ 生成的\Term{ideal}{理想}，便有
\[
K_\bullet(Y)=K_\bullet(S)\otimes_{K'(S)}(K'(X)/I).
\]
\SourcePage{527}
而由第 VI 讲 4.8，$K'(X)$ 是由各 $\xi_i-1$ 生成的 $K'(S)$-\Term{algebra}{代数}，其关系来自
\[
\gamma_t(E-n)=\prod_{i=1}^n\gamma_t(\xi_i-1).
\]
因此，$K'(X)/I=K'(S)/J$，其中 $J$ 是由 $\gamma^j(E-n)$（$j=1,\ldots,n$）生成的\Term{ideal}{理想}，也就是由 $E-n$ 生成的\Term{lambda-ideal}{$\lambda$-理想}。所以 $K_\bullet(Y)\simeq K_\bullet(S)/JK_\bullet(S)$。下面的\Term{lemma}{引理}完成证明。
\end{proof}

\begin{SGAresult}{\Term{lemma}{引理}}{4.2}
设 $S$ 是\Term{noetherian}{诺特}\Term{scheme}{概形}，$B$ 是某个\Term{reductive-group}{约化} $S$-群的\Term{parabolic-subgroup}{抛物子群}，$P$ 是 $B$ 的\Term{unipotent-radical}{幺幂根基} $U$ 下的\Term{torsor}{挠子}；则有\Term{isomorphism}{同构}
\[
K_\bullet(S)\xrightarrow{\sim}K_\bullet(P).
\]
\end{SGAresult}
\begin{proof}
事实上，存在 $U$ 的一列\Term{subgroup-scheme}{子群概形}
\[
U_0=U\supset U_1\supset U_2\supset\cdots\supset U_p=e,
\]
它们在 $B$ 中闭，且对每个 $i$，$U_i/U_{i+1}$ \Term{isomorphism}{同构}于一个 $S$-\Term{vector-group}{向量群}（SGA 3，第 XXVI 讲 2.1）；由于 $B$ 是\Term{noetherian}{诺特的}，此列是\Term{finite}{有限}的。因此，$P\to S$ \Term{resolution}{分解}为一连串\Term{vector-group}{向量群}下的\Term{torsor}{挠子}，由 3.4 即得结论。
\end{proof}

\begin{SGAresult}{推论}{4.3}
设 $S$ 是\Term{noetherian}{诺特}\Term{scheme}{概形}，$G=\mathrm{GL}(n)_S$。则\Term{canonical}{典范}\Term{homomorphism}{同态} $K_\bullet(S)\to K_\bullet(G)$ 是\Term{isomorphism}{同构}。
\end{SGAresult}
\begin{proof}
将 4.1 用于 $P=G$ 即得，因为这时 $\mathcal E=\mathcal O_S^n$。
\end{proof}

\SourcePage{528}
\begin{SGAresult}{推论}{4.4}
在 4.1 的假设下，再假设 $S$ \Term{regular}{正则}。则有\Term{isomorphism}{同构}
\[
K'(S)/J\xrightarrow{\sim}K'(P).
\]
\end{SGAresult}
\begin{proof}
若 $S$ \Term{regular}{正则}，则 $G$ 也\Term{regular}{正则}，$P$ 同样\Term{regular}{正则}。因此有 $K'(S)\simeq K_\bullet(S)$、$K'(P)\simeq K_\bullet(P)$，结果便由 4.1 得出。
\end{proof}
